\chapter{基于三维高斯的语义特征蒸馏与前景背景分离}

\section{引言}
在上一章中，系统基于多模态大模型和深度分层技术，完成了二维平面下的细粒度实例掩码提取与基础本构属性映射，并获得了带有空间位姿信息的初始点云。然而，对于后续的三维物理仿真（如物质点法）而言，纯粹的二维掩码和稀疏离散点云并不足以支撑连续、高保真的几何碰撞与材质运算。这就要求我们将从多重视图和二维基础模型中提取到的高维特征，持续且稳定地投影到具有明确几何体积的三维表征之中。
为此，本章以三维高斯溅射（3D Gaussian Splatting, 3DGS）为核心底层表示，设计了一套联合语言先验与视觉几何先验的特征蒸馏框架。本章首先阐述如何在标准高斯球参数中嵌入任意维度的特征向量，随后提出融合 CLIP 与 DINO 多尺度特征的双流蒸馏策略，并重点引入了具备自适应权重的 DINO 损失预热（Warmup）机制；最终结合空间梯度边界约束与 10 大类材质分类框架，实现背景剔除与前景复杂物理材质像素的精准解耦。这一架构为后续力学引擎赋予三维空间实体以真实的物理属性打下了坚实的场景表示基础。

\section{基于三维高斯溅射的物理语义场表征}
经典的三维高斯溅射技术凭借其显式的椭球体包围盒以及高效的可微光栅化渲染管线，在静态场景的新视角合成任务上取得了举世瞩目的成就。每个高斯原语不仅包含中心位置信息，还具备协方差矩阵（由缩放因子和旋转四元数定义）、透明度及由于球谐函数（SH）描述的颜色信息。
为了使高斯原语在表达视觉几何的同时能够承载力学仿真所需的语义与物质类别，我们扩展了高斯原语的基座属性集，为每个高斯球挂载两组独立的高维语义通道：

\textbf{视觉语义特征向量（Visual Semantic Feature）}：存储从 CLIP与 DINO 视觉基础模型中蒸馏提取的跨模态语义与细粒度几何特征，并进一步扩展为部件级（part-level）语义特征通道，以支持更细粒度的物体分区表征。

\textbf{材质语义标签向量（Material Semantic Label）}：存储 VLM（Qwen3-VL-32B）对场景各区域推理预测的材质类别标签，经语义标签蒸馏后同步映射至三维高斯空间，使每个高斯原语明确归属到具体的材质大类（如刚体、弹性体、织物等）。

在正向渲染过程中，与传统光栅化颜色的色彩累积相似，这两组高维语义通道同样遵循 $\alpha$-blending 的透射率积分公式，各自独立渲染并通过对应的解码器输出。通过赋予语义场可微的渲染能力，三维网格中的离散物理材质特征可以被顺滑地投影至任意二维视角，从而分别与 CLIP/DINO 视觉特征监督图以及 VLM 材质分类标签进行损失回传，实现三维高斯语义场在多维视觉语义与物理材质分类两个层面上的精准对齐。


\section{二维语义特征预提取：多层级视觉监督信号生成}

在将语义特征蒸馏至高斯场之前，系统需先从原始多视角图像中提取高质量的多层级语义特征作为监督目标。该预处理阶段由 \texttt{compute\_obj\_part\_feature.py} 模块独立完成，为后续的 Feature-Splatting 训练提供三类互补的视觉特征信号。

\subsection{物体感知分割与特征提取}
传统 SAM 的自动网格提示模式在面对复杂场景时，常常将目标对象分解为大量碎片化的区域，无法为每个独立物体生成完整的掩码。为解决这一问题，系统引入了物体感知的 YOLO-MobileSAMv2 联合检测-分割管线：首先由 YOLOv8 检测器对每帧图像进行目标检测，输出各物体的包围盒（bounding box）与置信度；随后将包围盒作为提示（Prompt）输入 MobileSAMv2 轻量级分割模型，获得每个物体的精确实例掩码。这一"检测引导分割"策略从根本上解决了无提示 SAM 区域碎片化的问题，确保每个材质部件都能被完整覆盖。

\subsection{对象级与零件级 CLIP 特征}
CLIP 特征的提取分为两个层次：
\textbf{对象级特征（Object-level Features）}：将整帧图像送入 CLIP ViT-L/14 编码器（输入分辨率 336 像素），获取保留空间布局的 patch-level 特征图（维度 $C \times H_{\text{patch}} \times W_{\text{patch}}$）。随后利用 YOLO-MobileSAMv2 获得的实例掩码对特征图进行空间掩码平均池化，得到每个检测对象的全局语义特征向量。

\textbf{零件级特征（Part-level Features）}：当系统配置为 \texttt{clip\_part} 模式时，对 YOLO 检测出的每个包围盒区域单独裁剪，缩放到统一分辨率后独立编码，获得零件级的局部语义特征；最后将各包围盒特征按空间位置回贴至图像网格中，形成保留细粒度材质分区的零件级特征图。两种特征经维度级联后统一通过双线性插值降维至目标分辨率（128 维），在信息完整性与存储效率之间取得平衡。

\subsection{DINOv2 空间结构特征}
为了弥补 CLIP 对物体边界和局部纹理敏感度不足的缺陷，系统同步从 DINOv2 ViT-S/14 自监督模型中提取 patch token 特征（输入最长边 1120 像素）。DINOv2 通过自蒸馏预训练习得的注意力机制天然对物体轮廓与高频纹理具有响应，其 patch 特征图可为后续的语义一致性约束提供空间平滑性引导。

经此阶段，每帧图像生成三类特征文件——\texttt{clip\_feat\_obj.npy}、\texttt{clip\_feat\_part.npy} 与 \texttt{dino\_feat\_with\_bbox.npy}——共同构成 Feature-Splatting 训练阶段的多源语义监督信号。
\section{多尺度视觉大模型特征的双流蒸馏}
为了对扩展后的三维高斯语义场进行有效监督，系统摒弃了单一维度的特征指引，转而采用了结合跨模态语义（CLIP）与细粒度视觉几何（DINO）的双流特征蒸馏方案，并默认启用部件级（Part-level）特征蒸馏以增强物体内部的结构化材质分区能力。

\subsection{CLIP 与 DINO 的特征互补机制}
CLIP 作为以文本-图像对比学习为核心的视觉语言模型，赋予了系统从开集自然语言描述到视觉区块的强尺度响应，使渲染视角能够理解"软橡胶"、"金属"等高级物质概念。而 DINO 作为自监督视觉特征提取网络，其输出特征层天然对物体边界、局部纹理及轮廓极其敏感，能够规避 CLIP 在生成掩码时常见的"感受野外溢"缺陷。
通过将这两种基础模型的特征维度投影至降维子空间（例如 32 维），网络能够在控制显存开销的前提下，计算渲染特征图与目标基准图在余弦相似度上的误差（Cosine Distance Loss）。系统引入了可配置的 DINO 归一化损失权重（dino\_norm\_loss\_weight，默认 0.0），替代了此前硬编码的 0.1 比例，使用户可根据场景复杂度灵活调整 DINO 范数匹配项的介入强度。此外，特征导向的语义一致性约束（Feature-guided Semantic Consistency）新增了多尺度下采样机制（max\_size 默认 192 像素），当特征图尺寸超过该阈值时自动通过双线性插值降采样，在有效保障空间中连续物体材质平滑过渡的同时大幅降低显存与计算开销。

\subsection{自适应 DINO 损失预热（Warmup）策略}
在实际工程实验中，我们发现如果在训练初期便施加过强的 DINO 特征监督，极易导致高斯原语将部分容量浪费于死板拟合抽象高维向量上，从而破坏最初依赖 RGB 收敛确立的三维空间几何骨架。
为了解决这一时序收敛冲突，我们在特征蒸馏算子中设计了一套非线性预热（Warmup）机制。具体而言，系统设置了 20,000 次迭代的爬坡周期。在此阶段内，DINO 损失的权重并非恒定，而是以一个极其保守的初始值（0.2）启动，并随着迭代批次的推进平滑爬升至最高阈值（0.45）。这种基于训练周期的预热函数保证了三维高斯场在起步阶段能够优先巩固其几何密度和形变张量，待物体表壳基本定型后再逐渐施加边缘敏锐的特征压迫。该机制在提高模型整体峰值信噪比（PSNR）的同时，大幅度削减了前景与背景接壤处的“半透明絮状”伪影。

\section{面向动力学的前景解耦与 10 分类材质框架}
成功建立连续的高维视觉语义场后，为了对接后续的物理计算管线（如 MPM），系统还必须将 VLM 预测的材质类别信息显式地编码到三维空间，实现从连续语义特征到离散材质标签的精确定位。

\subsection{10 类基础物理材质体系映射}
在本系统中，我们摒弃了过于冗杂繁复且易导致分类器坍塌的 19 类细分标准，最终重构了一套包含 10 个核心大类的精简材质分类框架。这 10 个类别（包含刚体、弹性体、织物、流体边界等）严格对应于力学解算器所支持的本构模型。
在网络架构上，上述两组语义向量分别由独立的轻量级解码器处理：视觉语义特征向量通过 CLIP/DINO 解码器重建视觉基础模型的高维特征图，用于跨模态语义蒸馏与训练监督；而材质语义标签向量则经由专用的语义标签解码器（Semantic Label Decoder）输出针对这 10 个通道的离散化分类 logits，将第三章中 VLM 推理得到的材质类别从二维掩码空间蒸馏至每个三维高斯原语。当网络生成具体的二维材质掩码时，系统利用 \texttt{semantic\_target \% 10} 及 \texttt{semantic\_target // 10} 等编码逻辑，有效剔除无效分类目标（Ignored targets）和混合模糊区域。这两个解码器在训练过程中共享骨干网络但独立优化各自的损失函数，确保视觉语义感知与物理材质分类互不干扰。

\subsection{长尾分布及边界加权的约束优化}
在进行真实场景材质归类时，某些特殊类别往往仅占据极少的视度像素，形成了典型的实例长尾分布效应（Long-tailed distribution）。为此，系统在损失评估中引入了区域反频率权重（Inverse-frequency weights），对稀有材质进行自适应误差放大，避免网络在梯度下降时向占地极广的背景或主要载体塌陷。同时，以 $\gamma = 0.5$ 的类别平衡指数对交叉熵损失进行指数加权，并在语义特征解码器中引入独立于视觉蒸馏的\textbf{语义特征库（Separate Semantic Feature Bank）}，使语义分类损失的梯度仅回传至语义专用特征而不会污染 CLIP/DINO 的视觉特征空间。为进一步确保三维语义一致性，系统还引入了\textbf{逐高斯语义一致性损失（Gaussian Semantic Consistency Loss）}：利用光栅化器输出的像素-高斯索引，将二维语义标签反投影至每个高斯原语进行投票监督，迫使各高斯在三维空间中保持统一的材质归属。语义损失的整体调度采用精心设计的时延策略——像素级交叉熵在 20,000 次迭代后接入，逐高斯约束则延迟至 25,000 次迭代并以 5,000 次线性爬坡缓启——确保在视觉特征充分收敛后才施加语义层级的结构化约束。
此外，由于深度混叠容易导致物质边界处的物理参数重叠交界，算法专门增加了类别边界感知加权机制。在该约束下，网络不会对处于深度陡变或是特征极度不连贯处（被标记为 ignored pixels）的不可靠像素盲目信任，而是通过局部特征的余弦平滑（Similarity threshold）迫使两侧各自保持自身的空间特征纯净性。为进一步防止复杂场景训练过程中的显存溢出（OOM），系统引入了高斯点密度上限控制（densify\_max\_points，默认 5,000,000），当高斯原语数量达到该阈值时自动跳过致密化操作，保障大场景训练的内存安全性。最终，原本复杂交织的背景底图被有效遮蔽剥离，而前景物体则分化为边界锋利、材质种类绝对纯粹的三维力学操作域。

\section{本章小结}
本章系统性地阐明了从二维图片到全三维几何物理场的隐式构建与蒸馏解耦方法。通过在标准高斯原语上挂载视觉语义特征与材质语义标签两组独立的高维向量，并分别配置 CLIP/DINO 特征解码器与语义标签解码器进行双线蒸馏，本网络使得每一处几何点不仅具备视觉透射率，更在视觉语义与物理材质两个维度上“知悉”其所在的语义范畴。独创性的 YOLO-MobileSAMv2 检测引导分割、DINO 递增预热训练（0.2→0.45，20,000 iter）、独立语义特征库双线训练、逐高斯语义一致性损失（Gaussian Consistency Loss）、特征引导的 8 邻域语义平滑约束、时延语义损失调度以及高斯点密度上限控制等多项优化，完美攻克了高维特征嵌入导致几何骨架破碎、多任务特征冲突、三维语义不一致与显存溢出等难题。配合改良落地的 10 维动力学材质分类框架与复合 ID 等级解码机制，虚拟场景实现了从外貌复刻走向材质分离质变的决定性一步阶段，为下文梦幻物理（DreamPhysics）的视频运动扩散生成及物质点法仿真铺平了道路。
